\documentclass[conference]{IEEEtran}
\usepackage{cite}
\usepackage{amsmath,amssymb,amsfonts}
\usepackage{graphicx}
\usepackage{textcomp}
\usepackage[T1]{fontenc}
\usepackage{url}
\usepackage{booktabs}
\usepackage{multirow}

\begin{document}

\title{UnibaBot PDA: An Intelligent Agent for Automated Compliance Verification of Academic Development Plans Using RAG and Fine-Tuned LLMs}

\author{
\IEEEauthorblockN{Kevin Beltr\'{a}n Mart\'{i}nez}
\IEEEauthorblockA{Code: 2220221003\\
Systems Engineering\\
Faculty of Engineering\\
University of Ibagu\'{e}\\
Ibagu\'{e}, Colombia}
\and
\IEEEauthorblockN{Jeryleefth Lasso Motta}
\IEEEauthorblockA{Code: 2220231074\\
Systems Engineering\\
Faculty of Engineering\\
University of Ibagu\'{e}\\
Ibagu\'{e}, Colombia}
}

\maketitle

\begin{abstract}
Every semester, academic program offices at the University of Ibagu\'{e} have to manually review more than twenty Academic Development Plans (PDAs) to make sure they follow institutional guidelines, ABET student outcomes, and the competency mappings defined for each course. This process takes a long time, produces inconsistent results, and is easy to get wrong. We propose UnibaBot PDA, an intelligent agent that automates this review by combining Retrieval-Augmented Generation (RAG) with a fine-tuned large language model (LLM). The institutional rules---including ABET outcomes with their performance indicators, program competencies, and per-course mappings---are stored in a vector database so the system can look up the right criteria for each course when it needs them. A pre-trained LLM, adapted using parameter-efficient fine-tuning (LoRA), then reads each PDA section alongside the retrieved requirements and produces a structured compliance report. This paper presents the problem, reviews the current state of research in RAG, LLM fine-tuning, and automated compliance checking, and lays out the goals and research questions that guide the project.
\end{abstract}

\begin{IEEEkeywords}
large language models, retrieval-augmented generation, fine-tuning, LoRA, compliance verification, ABET, academic quality assurance, intelligent agents
\end{IEEEkeywords}

% ==============================================================
\section{Introduction}
% ==============================================================

Every semester, the offices of the different academic programs at the University of Ibagu\'{e} have to manually review more than twenty Academic Development Plans (PDAs) to check that each one follows the institutional guidelines and improvement directives for that period. Right now, this is done entirely by hand: faculty and staff read through each document, compare it against a set of rules, and write down any problems they find. With so many documents to review and no way to eliminate human subjectivity, mistakes slip through, different reviewers apply the criteria differently, and the whole process adds significant delays to the academic calendar. These problems affect both the quality of the review and the ability to prepare courses on time.

PDAs follow a semi-standardized template with sections like learning outcomes, assessment methods, course schedules, and improvement actions from the previous semester. The rules they must follow are not a single document but a structured set of connected requirements: seven ABET student outcomes (O1--O7) with specific performance indicators, three program-level competencies (C1--C3), twelve generic competencies, five Saber Pro components, six institutional dimensions, and a mapping that tells exactly which of these elements each course must cover. For example, the course ``Agentes Inteligentes'' (code 22A14) must develop competencies C1 and C2, generic competencies 1c and 1h, Saber Pro component SP5, and institutional dimension D4. A PDA for this course needs to show, in its written sections, that it actually addresses all of these requirements. Both the PDAs and the rules are written in Spanish, which adds a language processing challenge.

To tackle this problem, we propose UnibaBot PDA, an intelligent agent that can automatically analyze each PDA against the institutional requirements and generate a structured compliance report. The system takes a PDA as a PDF file and produces a detailed report showing which sections do not meet the requirements, what specific rules were violated, and what corrections are needed. The goal is to cut the review time from days to minutes, get rid of the inconsistency that comes with human review, and let program staff focus on more important academic decisions.

The system combines two techniques that work together. First, Retrieval-Augmented Generation (RAG) \cite{lewis2020rag} stores the competency definitions, ABET outcome descriptions, and performance indicators in a vector database and pulls up the right criteria for each PDA section when needed. The per-course mappings (which competencies apply to which course) are structured data that can be looked up directly, but the descriptions that explain what each competency actually means are written in natural language and benefit from semantic search. Second, a pre-trained LLM is specialized through parameter-efficient fine-tuning---specifically LoRA \cite{hu2022lora}---on examples of compliance analysis to get better at evaluating academic documents against the rules.

The remainder of this paper is organized as follows. Section~II reviews the state of the art across four areas: large language models, retrieval-augmented generation, parameter-efficient fine-tuning, and automated compliance checking. Section~III formulates the research goals and questions. Section~IV describes the proposed system architecture. Section~V outlines the planned evaluation methodology. Section~VI concludes with a discussion of expected contributions.

% ==============================================================
\section{State of the Art}
% ==============================================================

This section reviews the foundational work and recent advances that inform the design of UnibaBot PDA, organized into four areas: large language models, retrieval-augmented generation, parameter-efficient fine-tuning, and automated compliance verification.

% --------------------------------------------------------------
\subsection{Large Language Models}
% --------------------------------------------------------------

Modern large language models are built on the Transformer architecture, introduced by Vaswani et al. \cite{vaswani2017attention}. What makes the Transformer special is its self-attention mechanism: it lets the model look at every word in a sentence and figure out which other words are most relevant to it, no matter how far apart they are. Previous models (recurrent neural networks) had to read text one word at a time from left to right, which made them slow and bad at connecting ideas that were far apart in a text. The Transformer reads everything at once, which makes it faster to train and much better at understanding context. Today, practically every major language model is built on this architecture.

Two main families of models came out of the Transformer. The first is the generative family, best known through GPT. Radford et al. \cite{radford2018gpt1} showed that training a Transformer on a huge amount of text first and then fine-tuning it for a specific task produces very strong results. GPT-2 \cite{radford2019gpt2} took this further and showed that bigger models can even handle tasks they were never explicitly trained for---a capability called zero-shot learning. The second family is the understanding-focused one, led by BERT \cite{devlin2019bert}, which reads text in both directions at the same time (left-to-right and right-to-left). BERT is especially good at tasks like classifying text and extracting information, which is exactly what compliance checking requires: determining whether a section of a document actually meets a given requirement.

What made all of this practical for projects like ours was the release of open-weight models---models whose weights are publicly available. LLaMA \cite{touvron2023llama} showed that smaller, well-trained models (from 7 billion to 65 billion parameters) can perform as well as or better than much larger proprietary ones. Mistral 7B \cite{jiang2023mistral} pushed efficiency even further with architectural improvements. These models can run on free cloud GPUs like Google Colab, which is what we have available. For our system, we are looking at models in the 3B--7B parameter range, since they offer a good balance between capability and the hardware we can actually use.

% --------------------------------------------------------------
\subsection{Retrieval-Augmented Generation}
% --------------------------------------------------------------

A key limitation of language models is that they only know what was in their training data. If something was not included---or if it changed after training---the model has no way to access it. Retrieval-Augmented Generation (RAG), introduced by Lewis et al. \cite{lewis2020rag}, solves this by giving the model access to an external knowledge base. When the system gets a question, it first searches the knowledge base for relevant information, then feeds both the question and the retrieved context to the language model so it can generate an informed answer. This way, the model always works with the right information instead of relying only on what it memorized.

RAG has evolved a lot since it was first proposed. Gao et al. \cite{gao2024ragsurvey} describe three stages of this evolution: Naive RAG, which simply retrieves and reads; Advanced RAG, which improves the query before searching and re-ranks results afterward; and Modular RAG, which uses interchangeable pieces that can be swapped depending on the domain. The most important takeaway from their survey for our project is that the way documents are split into chunks before storing them in the database makes a big difference in the quality of the search results.

Fan et al. \cite{fan2024ragmeeting} found that the quality of the retrieval step is the single most important factor in how good the final output is. If the system fails to find the right passage---even when it is there in the database---the answer it generates will be wrong no matter how smart the model is. This matters a lot for our project: if the system cannot find the right ABET indicator or competency description for a given PDA section, its compliance assessment will be unreliable.

Jimeno Yepes et al. \cite{jimenoyepes2024chunking} studied how to split financial reports into chunks for RAG, a domain similar to ours since both deal with semi-structured documents that have section headers and references to rules. They found that splitting documents at natural section boundaries (instead of cutting at a fixed number of characters) improves search precision by 15--20\%. This is directly useful for us: each ABET outcome paired with its indicators forms a natural chunk that should be kept together.

Sun et al. \cite{sun2025compliance} built a compliance checking system directly on top of RAG, where regulatory clauses are retrieved and fed to the LLM for assessment. Their system performs better than approaches that only use the model's built-in knowledge, which confirms that pulling in external rules at evaluation time is the right approach for compliance tasks. Their architecture is very similar to what we are building for PDA verification.

In our system, the knowledge base contains the textual descriptions of ABET outcomes and their indicators, competency definitions, and Saber Pro component descriptions. The per-course mappings (which competencies apply to which course) are structured data stored in JSON format and accessed through direct lookup by course code. RAG is used specifically for the semantic matching step: given a free-text PDA section, the system retrieves the most relevant competency and outcome descriptions to determine whether the section adequately addresses the required criteria.

% --------------------------------------------------------------
\subsection{Parameter-Efficient Fine-Tuning}
% --------------------------------------------------------------

Pre-trained language models are good at many things, but they usually get better when trained further on data from a specific domain. The problem is that retraining all the parameters of a large model takes a lot of computing power---a 7-billion-parameter model needs over 80GB of GPU memory for full fine-tuning, which is far beyond what free cloud platforms offer.

Parameter-efficient fine-tuning methods get around this by only training a small number of extra parameters while leaving the original model untouched. LoRA (Low-Rank Adaptation), proposed by Hu et al. \cite{hu2022lora}, is the most popular one. The idea is that the adjustments a model needs to learn a new task can be captured by small matrices added on top of the existing weights. In practice, this means training less than 1\% of the total parameters while still getting results close to full fine-tuning. Once training is done, the new weights can be folded back into the original model, so there is no extra cost when running it.

QLoRA, introduced by Dettmers et al. \cite{dettmers2023qlora}, takes LoRA a step further by first compressing the base model to 4-bit precision before adding the adapters. This cuts memory use even more---a 7B-parameter model that would normally need 16GB fits in roughly 6GB---which makes fine-tuning possible on free Google Colab T4 GPUs (16GB of VRAM). Despite this aggressive compression, Dettmers et al. showed that the quality loss is minimal: QLoRA-trained models reached 99.3\% of the performance of models trained at full precision.

Instruction tuning, demonstrated by Wei et al. with FLAN \cite{wei2022flan}, showed that training a model on examples formatted as instruction-response pairs enables it to follow new instructions it has never seen before. This is the format we plan to use for our training data: each example is an instruction like ``Evaluate whether this PDA section addresses competency C1 and ABET indicator 1.1'' paired with a response that says whether it complies and what specific issues exist.

Taori et al. \cite{taori2023alpaca} showed that these training datasets can be created cheaply by using a more powerful model to generate them. They used GPT-3.5 to produce 52,000 synthetic instruction-response pairs and then fine-tuned a smaller LLaMA 7B model on that data, getting competitive results for under \$600. This self-instruct approach \cite{wang2022selfinstruct} is what we plan to do: use a capable model to generate synthetic compliance-analysis examples from real PDA sections combined with the competency mappings, creating training data without having to write it all by hand.

% --------------------------------------------------------------
\subsection{Automated Compliance Checking}
% --------------------------------------------------------------

Automated compliance checking has been explored in several other fields, but almost no one has applied it to academic documents.

Amaral Cejas et al. \cite{amaral2023gdpr} used NLP to check data processing agreements against GDPR requirements. Their system pulls obligations out of the regulation, maps them to sections of the document, and classifies each section as compliant or not. Even though their domain is different, the approach---break the rules into pieces and check the document against each one---is the same pattern we use for PDA verification.

Zheng et al. \cite{zheng2025construction} review construction contracts using LLMs enhanced with a knowledge graph of regulations. Their main finding is that combining structured knowledge with the LLM's language abilities produces better results than using either one alone. This directly supports our approach of mixing structured JSON lookups with LLM-based text analysis.

Savelka et al. \cite{savelka2023legal} found that LLMs can spot regulatory obligations in legal texts with high accuracy even without any domain-specific training (zero-shot). This tells us that modern LLMs already have a solid starting point for compliance tasks, and fine-tuning can push them even further.

Sun et al. \cite{sun2025docagent} propose DocAgent, a system that breaks complex document analysis into smaller sub-tasks handled by specialized components. This modular approach is relevant to PDA review, where different sections (learning outcomes, assessment methods, improvement actions) may need different analysis strategies.

The biggest gap we found in the literature is that nobody has built a system for academic compliance verification, especially not for ABET-accredited programs or documents written in Spanish. Our work fills that gap by combining structured competency data with RAG-enhanced LLM analysis, applied specifically to university PDA review.

% ==============================================================
\section{Research Goals and Questions}
% ==============================================================

\subsection{General Objective}

To design and implement an intelligent agent that automatically verifies the compliance of Academic Development Plans (PDAs) against the institutional competency framework---including ABET student outcomes, program-specific competencies, generic competencies, Saber Pro components, and institutional dimensions---generating structured reports that identify non-compliant sections, the specific criteria violated, and the corrections required.

\subsection{Specific Objective}

\begin{enumerate}
    \item To evaluate the system's accuracy and consistency against human expert reviews on a representative set of PDAs from the Systems Engineering program.
\end{enumerate}

\subsection{Research Questions}

\begin{enumerate}
    \item \textbf{RQ1:} To what extent can a RAG-augmented, fine-tuned LLM accurately identify non-compliant sections in PDAs compared to human expert reviewers?
    \item \textbf{RQ2:} How does the combination of structured competency lookup and RAG retrieval with domain-specific fine-tuning compare to using either technique in isolation for compliance verification accuracy?
\end{enumerate}

% ==============================================================
\section{Proposed System Architecture}
% ==============================================================

UnibaBot PDA operates as a pipeline with four stages, illustrated in Fig.~\ref{fig:architecture_placeholder}.

\textbf{Stage 1: Document Ingestion.} The system receives a PDA in PDF format, extracts the text, and segments it into sections based on the document's heading structure (learning outcomes, assessment methods, course schedule, improvement actions, etc.).

\textbf{Stage 2: Course Identification and Requirement Lookup.} The system identifies the course code from the PDA (e.g., 22A14 for ``Agentes Inteligentes'') and performs a direct lookup in the structured competency mapping to retrieve the list of required competencies, ABET indicators, generic competencies, Saber Pro components, and institutional dimensions for that course. This step does not require semantic search---it is a deterministic lookup against the JSON-encoded course-competency mapping.

\textbf{Stage 3: Semantic Retrieval and Compliance Analysis.} For each required competency or ABET indicator, the system uses RAG to retrieve the full textual description from the vector store (e.g., for ABET indicator 1.1: ``Identify the problems and applicable theories and concepts''). It then passes each (PDA section, requirement description) pair to the LLM, which determines whether the PDA section adequately addresses the requirement and, if not, explains what is missing or incorrect.

\textbf{Stage 4: Report Generation.} The per-requirement assessments are aggregated into a complete compliance report listing: requirements that are fully addressed, requirements that are partially or not addressed (with specific comments), and any PDA content that does not correspond to the expected format.

The system is designed to be consistent: the same PDA evaluated against the same rules will always produce the same result. It processes documents one at a time in sequence and delivers the full set of compliance reports at the end.

For the fine-tuning part, we plan to use LoRA \cite{hu2022lora} on a model in the 3B--7B range (such as Llama 3.2 3B \cite{touvron2023llama}), chosen because it runs on Google Colab. If memory is too tight, we will switch to QLoRA \cite{dettmers2023qlora} with 4-bit compression. The training data will be instruction-response pairs built from real PDA sections and their corresponding competency requirements, using the self-instruct method \cite{wang2022selfinstruct}.

\begin{figure}[ht]
\centering
\fbox{\parbox{0.9\columnwidth}{\centering\vspace{1em}\textit{Architecture diagram to be included in a future revision.}\vspace{1em}}}
\caption{Overview of the UnibaBot PDA pipeline. Stage 1 extracts and segments the PDA. Stage 2 performs a deterministic lookup of course requirements. Stage 3 uses RAG to retrieve competency descriptions and the LLM to assess compliance. Stage 4 generates the final report.}
\label{fig:architecture_placeholder}
\end{figure}

% ==============================================================
\section{Evaluation Methodology}
% ==============================================================

The system will be evaluated along three dimensions.

\textbf{Correctness} measures the accuracy of compliance assessments. For each PDA section and each required competency, the system's judgment (compliant, partially compliant, or non-compliant) will be compared against human expert annotations. We will report precision, recall, and F1 score for the detection of non-compliant sections.

\textbf{Granularity} evaluates whether the system correctly identifies not only which sections fail but also the specific requirement violated and what correction is needed. The generated comments will be compared against reference comments written by human reviewers using ROUGE-L \cite{lin2004rouge} for surface-level overlap and BERTScore \cite{zhang2020bertscore} for semantic similarity.

\textbf{Consistency} measures whether the system produces the same assessment for the same PDA across multiple runs, and whether it applies criteria uniformly across different PDAs---a key advantage over human review, where inter-reviewer agreement is often low.

Additionally, a comparison study will evaluate three configurations: (a) a base LLM with structured lookup only (no RAG, no fine-tuning), (b) the base LLM with RAG retrieval (no fine-tuning), and (c) the full system with RAG and fine-tuning. This comparison directly addresses RQ2 by isolating the contribution of each component.

% ==============================================================
\section{Conclusion}
% ==============================================================

This paper has presented UnibaBot PDA, an intelligent agent for automating the compliance review of Academic Development Plans at the University of Ibagu\'{e}. The system combines direct lookup of per-course competency mappings, RAG-based retrieval of ABET outcome and competency descriptions, and a fine-tuned LLM for analyzing the actual text of each PDA. Our review of the literature shows that while RAG and fine-tuning have been used successfully for compliance checking in legal and regulatory fields, nobody has applied them to academic quality assurance---especially not for ABET-accredited programs or documents in Spanish. The architecture, research questions, and evaluation plan we have laid out provide a clear path toward turning the PDA review process from a manual, multi-day task into an automated one that is both faster and more consistent.

% ==============================================================
% REFERENCES
% ==============================================================
\begin{thebibliography}{00}

\bibitem{lewis2020rag}
P. Lewis et al., ``Retrieval-Augmented Generation for Knowledge-Intensive NLP Tasks,'' \textit{Adv. Neural Inf. Process. Syst.}, vol. 33, pp. 9459--9474, 2020, arXiv:2005.11401.

\bibitem{gao2024ragsurvey}
Y. Gao et al., ``Retrieval-Augmented Generation for Large Language Models: A Survey,'' 2024, arXiv:2312.10997.

\bibitem{fan2024ragmeeting}
W. Fan et al., ``A Survey on RAG Meeting LLMs: Towards Retrieval-Augmented Large Language Models,'' in \textit{Proc. 30th ACM SIGKDD Conf. Knowl. Discovery Data Mining (KDD)}, 2024, pp. 6491--6501.

\bibitem{sun2025compliance}
J. Sun, Z. Luo, and Y. Li, ``A Compliance Checking Framework Based on Retrieval Augmented Generation,'' in \textit{Proc. 31st Int. Conf. Comput. Linguistics (COLING)}, 2025, pp. 2603--2615.

\bibitem{jimenoyepes2024chunking}
A. Jimeno Yepes, Y. You, J. Milczek, S. Laverde, and R. Li, ``Financial Report Chunking for Effective Retrieval Augmented Generation,'' 2024, arXiv:2402.05131.

\bibitem{hu2022lora}
E. J. Hu et al., ``LoRA: Low-Rank Adaptation of Large Language Models,'' in \textit{Proc. Int. Conf. Learn. Representations (ICLR)}, 2022, arXiv:2106.09685.

\bibitem{dettmers2023qlora}
T. Dettmers, A. Pagnoni, A. Holtzman, and L. Zettlemoyer, ``QLoRA: Efficient Finetuning of Quantized LLMs,'' \textit{Adv. Neural Inf. Process. Syst.}, vol. 36, 2023, arXiv:2305.14314.

\bibitem{touvron2023llama}
H. Touvron et al., ``LLaMA: Open and Efficient Foundation Language Models,'' 2023, arXiv:2302.13971.

\bibitem{jiang2023mistral}
A. Q. Jiang et al., ``Mistral 7B,'' 2023, arXiv:2310.06825.

\bibitem{vaswani2017attention}
A. Vaswani, N. Shazeer, N. Parmar, J. Uszkoreit, L. Jones, A. N. Gomez, \L. Kaiser, and I. Polosukhin, ``Attention Is All You Need,'' in \textit{Proc. Adv. Neural Inf. Process. Syst. (NeurIPS)}, 2017, arXiv:1706.03762.

\bibitem{radford2018gpt1}
A. Radford, K. Narasimhan, T. Salimans, and I. Sutskever, ``Improving Language Understanding by Generative Pre-Training,'' OpenAI Technical Report, 2018.

\bibitem{radford2019gpt2}
A. Radford, J. Wu, R. Child, D. Luan, D. Amodei, and I. Sutskever, ``Language Models are Unsupervised Multitask Learners,'' OpenAI Technical Report, 2019.

\bibitem{devlin2019bert}
J. Devlin, M.-W. Chang, K. Lee, and K. Toutanova, ``BERT: Pre-Training of Deep Bidirectional Transformers for Language Understanding,'' in \textit{Proc. NAACL}, 2019, arXiv:1810.04805.

\bibitem{wei2022flan}
J. Wei, M. Bosma, V. Y. Zhao, K. Guu, et al., ``Finetuned Language Models Are Zero-Shot Learners,'' in \textit{Proc. Int. Conf. Learn. Representations (ICLR)}, 2022, arXiv:2109.01652.

\bibitem{taori2023alpaca}
R. Taori, I. Gulrajani, T. Zhang, Y. Dubois, et al., ``Alpaca: A Strong, Replicable Instruction-Following Model,'' Stanford CRFM, 2023.

\bibitem{wang2022selfinstruct}
Y. Wang et al., ``Self-Instruct: Aligning Language Models with Self-Generated Instructions,'' 2022, arXiv:2212.10560.

\bibitem{amaral2023gdpr}
O. Amaral Cejas, M. I. Azeem, S. Abualhaija, and L. C. Briand, ``NLP-Based Automated Compliance Checking of Data Processing Agreements Against GDPR,'' \textit{IEEE Trans. Softw. Eng.}, vol. 49, no. 9, pp. 4282--4303, 2023.

\bibitem{zheng2025construction}
C. Zheng, S. Wong, X. Su, Y. Tang, A. Nawaz, and M. Kassem, ``Automating Construction Contract Review Using Knowledge Graph-Enhanced Large Language Models,'' \textit{Autom. Constr.}, vol. 175, art. no. 106179, 2025.

\bibitem{savelka2023legal}
J. Savelka, K. D. Ashley, M. A. Gray, H. Westermann, and H. Xu, ``The Unreasonable Effectiveness of Large Language Models in Zero-Shot Semantic Annotation of Legal Texts,'' \textit{Front. Artif. Intell.}, vol. 6, art. no. 1279794, 2023.

\bibitem{sun2025docagent}
L. Sun, L. He, S. Jia, Y. He, and C. You, ``DocAgent: An Agentic Framework for Multi-Modal Long-Context Document Understanding,'' in \textit{Proc. Conf. Empirical Methods Natural Lang. Process. (EMNLP)}, 2025, pp. 17701--17716.

\bibitem{yao2023react}
S. Yao et al., ``ReAct: Synergizing Reasoning and Acting in Language Models,'' in \textit{Proc. Int. Conf. Learn. Representations (ICLR)}, 2023, arXiv:2210.03629.

\bibitem{lin2004rouge}
C.-Y. Lin, ``ROUGE: A Package for Automatic Evaluation of Summaries,'' in \textit{Text Summarization Branches Out}, 2004, pp. 74--81.

\bibitem{zhang2020bertscore}
T. Zhang, V. Kishore, F. Wu, K. Q. Weinberger, and Y. Artzi, ``BERTScore: Evaluating Text Generation with BERT,'' in \textit{Proc. Int. Conf. Learn. Representations (ICLR)}, 2020, arXiv:1904.09675.

\end{thebibliography}

\end{document}
